
base fidèlement plat quasi-compact. Si @@M0@@ est régulière, alors l'Idéal @@M1@@ qui définit @@M2@@ est de type fini, et @@M3@@ est localement libre de type fini, car ces propriétés se descendent par morphismes fidèlement plats quasi-compacts. Si @@M4@@ est une suite de sections de @@M5@@ relevant une base de @@M6@@, le complexe de Koszul associé aux images inverses des @@M7@@ est acyclique en degrés @@M8@@ d'après 1.4; comme c'est l'image inverse du complexe de Koszul associé aux @@M9@@, ce dernier est acyclique en degrés @@M10@@, et @@M11@@ est régulière.

\textbf{Définition 1.6.} \uline{Soit} @@M12@@ \uline{une immersion définie dans un ouvert} @@M13@@ \uline{de} @@M14@@ \uline{par un Idéal} @@M15@@ \uline{de} @@M16@@; \uline{le} @@M17@@-\uline{Module} @@M18@@ (\uline{indépendant de} @@M19@@) \uline{est appelé} faisceau conormal de @@M20@@ dans @@M21@@, \uline{et noté} @@M22@@. \uline{Si} @@M23@@ \uline{est régulière}, @@M24@@ \uline{est localement libre de type fini}, \uline{et son rang s'appelle} codimension de @@M25@@ dans @@M26@@, \uline{ou encore} \uline{codimension de} @@M27@@.

Si on a deux immersions fermées : @@M28@@, les faisceaux conormaux forment une suite exacte:
@@M29@@
Soient en effet @@M30@@ l'idéal de @@M31@@ dans @@M32@@, @@M33@@ l'idéal de @@M34@@ dans @@M35@@, @@M36@@ l'idéal de @@M37@@ dans @@M38@@. On a la suite exacte:
@@M39@@
d'où par tensorisation par @@M40@@, la suite exacte:
@@M41@@
en tensorisant par @@M42@@, on obtient la suite exacte:
@@M43@@
qui est la suite cherchée.

